\documentclass[conference,compsoc]{IEEEtran}
%\documentclass[conference, letter]{IEEEtran}
%\documentclass[draft,onecolumn,12pt]{IEEEtran}
%\documentclass[draft,onecolumn,12pt]{IEEEtran}
%\documentclass[journal]{IEEEtran}
%\documentclass[draftclsnofoot, onecolumn]{IEEEtran}
%\usepackage{graphicx,type1cm,eso-pic,color}

\usepackage{graphicx}
\usepackage{cite}
\usepackage{relsize}
\usepackage{xcolor}
\usepackage{colortbl}
\usepackage[cmex10]{amsmath}
\usepackage{amssymb}
\usepackage{amsmath}
\usepackage{booktabs}
\usepackage{multirow}
\usepackage{dcolumn}
\usepackage{amsthm}
\usepackage{enumerate}
\usepackage{algorithmicx}
\usepackage{array}
\usepackage{eqparbox}
\usepackage{verbatim}
\usepackage{bm}
%\usepackage[center]{caption}
\usepackage[font=footnotesize,justification=centering,singlelinecheck=false]{caption}
%\usepackage[caption=false,font=footnotesize]{subfig} % only if needed
\usepackage[utf8]{inputenc}
\usepackage{url}
\usepackage{enumitem}
\usepackage[]{footmisc}
%\usepackage{textcomp}
\usepackage{algorithm}% http://ctan.org/pkg/algorithm
\usepackage{algpseudocode}% http://ctan.org/pkg/algorithmicx
\usepackage{soul}

\usepackage{pgf}
\usepackage{siunitx}
% --- In preamble ---
\usepackage{tikz}
\usetikzlibrary{positioning,arrows.meta,shapes.geometric}
%\usepackage{glossaries}
%\makeglossaries % generates the acronym list
%\newglossaryentry{acm}{name={ACM},description={This is an elite glossary}}
%\newacronym{acm}{ACM}{Association for Computing Machinery}
\newtheorem{lemma}{Lemma}
\newtheorem{theorem}{Theorem}
\newtheorem{problem}{Problem}
\newtheorem{proposition}{Proposition}
\newtheorem{strategy}{Strategy}
\newtheorem{remark}{Remark}
\newtheorem{scenario}{Scenario}

\newcommand{\ins}[1]{\textcolor{blue}{#1}}
\newcommand\ddfrac[2]{\frac{\displaystyle #1}{\displaystyle #2}}
\renewcommand{\qed}{$\blacksquare$}
\DeclareMathOperator*{\argmin}{arg\,min}

% ================= COLORS =================
% Matte OrRd-style (white -> muted red), tuned to match figures
\definecolor{asrbase}{RGB}{140,13,18}    % matte deep red (high ASR)
\definecolor{asrdef}{RGB}{0,150,80}      % Green LLMaaD
\definecolor{asrimp}{RGB}{0,120,60}      % Darker green for reduction
% ================= CELL MACROS =================
% Baseline ASR (0-1 scale)  [white -> matte red]
\providecommand{\basecell}[1]{%
    \pgfmathtruncatemacro{\pct}{round(100*#1)}%
    \edef\colorspec{asrbase!\pct!white}%
    \expandafter\cellcolor\expandafter{\colorspec}%
    \pgfmathsetmacro{\v}{#1}%
    \ifdim \v pt > 0.70pt
    \textcolor{white}{\bfseries\num[round-mode=places,round-precision=3]{#1}}%
    \else
    \textcolor{black}{\bfseries\num[round-mode=places,round-precision=3]{#1}}%
    \fi
}
% LLMaaD ASR (0-1 scale)
\providecommand{\defcell}[1]{%
    \pgfmathtruncatemacro{\pct}{round(100*#1)}%
    \edef\colorspec{asrdef!\pct!white}%
    \expandafter\cellcolor\expandafter{\colorspec}%
    \num[round-mode=places,round-precision=3]{#1}%
}
% Percentage reduction (0-100%)
\providecommand{\impcell}[1]{%
    \pgfmathtruncatemacro{\pct}{round(#1)}%
    \edef\colorspec{asrimp!\pct!white}%
    \expandafter\cellcolor\expandafter{\colorspec}%
    \num[round-mode=places,round-precision=1]{#1}\%%
}
% Generic red cell (value in [min,max]) using the same matte base
\providecommand{\redcell}[3]{%
    \pgfmathsetmacro{\pctraw}{100*(#1-#2)/(#3-#2)}%
    \pgfmathtruncatemacro{\pct}{round(max(0,min(100,\pctraw)))}%
    \edef\colorspec{asrbase!\pct!white}%
    \expandafter\cellcolor\expandafter{\colorspec}%
    \num[round-mode=places,round-precision=3]{#1}%
}

\IEEEoverridecommandlockouts

\def \it {\textit}
\def \tbf {\textbf}
\def \mtc {\mathcal}
\def \balpha {\bm{\alpha}}
\def \bbeta {\bm{\beta}}
\def \bpi {\bm{\pi}}
\def \ba {\begin{array}}
\def \ea {\end{array}}
\def \benu {\begin{enumerate}}
\def \eenu {\end{enumerate}}
\def \bdes {\begin{description}}
\def \edes {\end{description}}
\def \itm {\itemize}
\def \bitem {\begin{itemize}}
\def \eitem {\end{itemize}}
\def \rz {\right .}
\def \ew {\mbox{elsewhere.}}
\def \ow {\mbox{otherwise.}}
\def \bfl {\begin{flushleft}}
\def \efl {\end{flushleft}}
\def \bfr {\begin{flushright}}
\def \efr {\end{flushright}}
\def \ni {\noindent}
\def \beq {\begin{equation}}
\def \eeq {\end{equation}}
\def \bqa {\begin{eqnarray}}
\def \eqa {\end{eqnarray}}
\def \bqa* {\begin{eqnarray*}}
\def \eqa* {\end{eqnarray*}}
\def \noi {\noindent}
\def \bal {\begin{align}}
\def \eal {\end{align}}

\def \tclr{\textcolor{red}}
\def \reza{\textcolor{magenta}}
\def \todo{\textcolor{orange}}
\newcommand{\mycmt}[1]{{\color{orange}\footnotesize[x]}}
\begin{document}

    % \title{Never Refuse, But Confuse: Enhancing LLM Defense Against Iterative Prompt Injection \& Jailbreak Attacks}
    %\title{\textit{Defensive Misdirection:} Enhancing LLM Security Against Prompt Fuzzing}
    %\title{LLMaaD: Misdirection as a Defense Against Agentic LLM Attacks}
%    \title{\emph{LLMaaD:} Misdirection as a Defense Against Model-Guided Automated LLM Attacks}
%    \title{Leveraging Misdirection in Defense Against Model-Guided Automated Attacks in Agentic AI}
    \title{Misdirection as a Defense Against Model-Guided Automated Attacks on Agentic AI Systems}

    \author{
        Reza Soosahabi, Vivek Namsani \thanks{Corresponding author: Reza Soosahabi (reza.soosahabi@keysight.com).}\\
        Application \& Threat Intelligence Research Center\\
        Keysight Technologies Inc., USA \\
    }


    \maketitle
    \thispagestyle{plain}\pagestyle{plain}


\begin{abstract}
 As generative AI is integrated into agentic applications, defenses increasingly combine model-level safeguards and external security controllers to monitor and mediate agent interactions.
 At the same time, attackers are adopting model-guided automation to scale multi-step probing, semantic prompt refinement, and response evaluation beyond manual effort.
 We analyze this setting probabilistically and derive asymptotic bounds on attacker success rate (ASR) as a function of response evaluation error rates.
 The analysis shows that conventional detect-and-block defense strategies can allow ASR to approach one as the attacker query budget grows.
 To address this limitation, we propose a detect-and-misdirect defense strategy that replaces predictable refusal responses with controlled, non-operational responses designed to corrupt the attacker’s automated response evaluation process through misdirection-induced false-positives.
We show that such false-positives result in a bounded asymptotic ASR.
 We instantiate this idea through Contextual Misdirection via Progressive Engagement (CMPE), a lightweight conversational misdirection method, and evaluate it on jailbreak benchmarks against multiple automated judge models.
 CMPE substantially increases attacker-side false-positives and reduces estimated ASR by up to two orders of magnitude across simulated attacker-defender judge configurations.
 We further evaluate CMPE against representative model-guided attack frameworks, PAIR and GPTFuzz, where it substantially reduces verified attack success rate and induces premature termination.
\end{abstract}

 \section{Introduction}
 Generative AI systems are increasingly deployed in agentic environments, where they can autonomously interact with external tools, retrieve untrusted content, and communicate with other agents.
 This expands the attack surface for jailbreak and prompt-injection attempts, especially when attackers can leverage automation to probe a target system and refine prompts over multiple interactions.
 While recent alignment methods improve resistance to isolated attempts, emerging model-guided automated attack frameworks use AI/ML models to generate and evaluate candidate prompts at scale, significantly amplifying attack efficiency \cite{gcg2023, llmfuzzer2024, pair2023}.
 A key characteristic of these attacks is the use of an automated judge component to evaluate the target model responses and provide feedback information for prompt refinement.
 
 \begin{figure}[t]
     \centering
     \includegraphics[width=\linewidth]{figs/cmpe-example.png}
     \caption{Example of contextual misdirection response generated by CMPE method in Algorithm \ref{algo:cmpe} against an automated jailbreak request.}
     \label{fig:cmpe}
 \end{figure}
 
In parallel, most deployed defenses rely on detection-based mechanisms followed by blocking or refusal, where malicious inputs or outputs are identified and suppressed through filtering or policy enforcement \cite{surveyusenix2024, wei2023jailbreak}.
Although effective against isolated attempts, such defenses produce structured and predictable outputs that can be exploited by attackers in automated search processes.
Modern model-guided attacks can use this feedback to iteratively improve prompts, either through optimization-based procedures in white-box settings \cite{gcg2023} or via fuzzing-style heuristics in black-box settings \cite{pair2023, llmfuzzer2024, autodan2023}.
In both cases, the defense response itself becomes part of the attack intelligence.
In Section~\ref{sec:theory}, we explore a fundamental limitation of detect-and-block defenses against model-guided automated attacks, arising from this feedback structure.
We formalize the attack loop using a probabilistic framework that models the interaction between the target system, including the defense mechanism, and the attacker’s judge.
Within this framework, we evaluate the attacker success rate (ASR) under a constrained verification setting and show that, under a broad class of search processes, the attacker success probability against such defenses can approach one as the attacker increases its query budget.
These findings are consistent with recent results demonstrating the existence of alignment failures that can be discovered through repeated querying \cite{su2024mission_impossible}. 
  
To address this limitation, we shift our focus from further improving detection accuracy to degrading the attacker’s feedback quality.
In Section~\ref{sec:propanalysis}, we propose a complementary defense strategy termed \textit{detect-and-misdirect}.
Instead of returning predictable refusal responses after detecting malicious behavior, the defense mechanism can produce in-context but misleading responses designed to interfere with the attacker’s evaluation process.
The key intuition is that model-guided attacks rely on the ability of automated judges to reliably distinguish successful adversarial outputs from failures.
By introducing semantically plausible but non-operational responses that appear successful to the \tclr{attacker's judge}, the defender can degrade the attacker's positive predictive value and diminish the efficiency of iterative search.
Section~\ref{sec:propanalysis} also formalizes this effect by introducing misdirection-induced \tclr{false-positives} in the attacker's evaluation process and shows that this leads to a bounded ASR even as the number of attack iterations grows.

At a practical level, we instantiate this strategy through a lightweight conversational mechanism called \textit{Contextual Misdirection via Progressive Engagement (CMPE)}, described in Section~\ref{sec:cmpe}.
CMPE replaces \tclr{predictable} jailbreak refusal text with responses that combine positive-intent framing, safe contextual expansion, and follow-up engagement.
While related to techniques studied in computational persuasion \cite{computational_persuasion}, the objective here is not to influence human users, but to disrupt automated adversarial evaluation.
This distinction is critical, as CMPE leverages known limitations of LLM-based judges, which often rely on heuristic cues such as tone and structure rather than strict semantic correctness \cite{zheng2023judging}.

We evaluate the proposed strategy in two complementary ways.
Section~\ref{sec:num} first estimates judge error rates using malicious and CMPE-generated misdirection responses over jailbreak benchmarks, and uses these estimates to compute ASR across simulated attacker-defender judge configurations.
These results show that CMPE increases attacker-side false-positive rates and reduces estimated maximum ASR by up to two orders of magnitude compared with the detect-and-block baseline.
More importantly, Section~\ref{sec:e2e_eval} evaluates CMPE in end-to-end attack runs against PAIR and GPTFuzz, where a vulnerable victim model is protected by a defense judge.
In these experiments, CMPE substantially reduces verified attack success and can cause automated attack frameworks to terminate prematurely on misdirection responses.

Section~\ref{sec:concdisc} concludes the paper by discussing deployment implications and future directions for integrating misdirection into agentic security architectures.
Overall, the paper contributes a probabilistic analysis of model-guided attack dynamics, a detect-and-misdirect defense strategy that bounds attacker success by degrading automated evaluation, and a practical CMPE instantiation validated through both judge-based ASR simulations and end-to-end attack-framework experiments.
These results support a defense perspective in which limiting the quality of attacker feedback can be as important as blocking malicious outputs in autonomous AI systems.

\section{Analysis of Model-guided Attack Dynamics}
\label{sec:theory}

There has been significant progress in securing generative AI models against prompt injection attacks as they are increasingly deployed in agentic applications with greater autonomy. Defenses have been integrated both into the models themselves and into external security layers that monitor agent interactions.
Models with reasoning capabilities such as chain-of-thought (CoT) mechanism have adopted strategies such as deliberative alignment (DA) that prevent user-induced deviations from system instructions or safety policies \cite{da2024}.
There are also structured query methods such as StruQ \cite{chen2025struq} for securing models by explicitly separating trusted instructions from untrusted data at the interface level, preventing prompt-injected instructions from being interpreted as executable commands.
These mechanisms explicitly enforce prioritization between trusted system instructions and untrusted external inputs, making direct prompt overrides and input semantic manipulations substantially less effective.
Thus, many emerging attack techniques involve multiple interactions with the model to probe and bypass defense mechanisms at deeper layers \cite{surveyusenix2024,manyshot2024}, which are time-prohibitive without some level of automation.
 This has motivated model-guided automated attack frameworks where AI/ML models are used to accelerate the evaluation of the sheer number of victim model responses and the generation of adversarial prompts with minimal human supervision.

\begin{figure}[t]
    \centering
    \includegraphics[width=\linewidth]{figs/bb-llm-fuzzing.png}
    \caption{Model-guided attacks \& the current detect-and-block defense strategy}
    \label{fig:fuzzing}
\end{figure}


Figure \ref{fig:fuzzing} depicts the components involved in a model-guided attack scenario: victim model, defense system, and attacker.
 The defense mechanism can be the reinforced behavior of the victim model, or a designated AI firewall system enforcing security policies on the victim system.
 Many current defense mechanisms follow a \textit{detect-and-block} strategy, where suspected malicious inputs are identified and handled through refusal messages, filtering, or interaction termination.
 They often scan for malicious patterns in both inbound and outbound content and explicitly react to them which makes their method of action known to external entities.

The attacker is also equipped with models to evaluate (judge) the reaction of the victim system to attack prompts, whose outputs serve as feedback information to the attack prompt generation mechanism.
 The algorithm \ref{algo:attack} illustrates a common workflow of model-guided attacks analyzed in this work.
 In each generation cycle, a set of $N$ prompts are created by a model or an algorithm based on a heuristic template, and then tested against the victim model, whose responses are evaluated by a judge model in the form of a score or a distance from the desired response.
 Prompts with high judge scores are placed in a candidate set.
 Then the attacker may manually verify up to $K$, often a small number, of prompt-response pairs in the candidate set  at the end of each cycle to avoid false-positives and ensure their potency and reproducibility, which is a time-consuming task.
 Since exhaustive verification of all $N$ responses is impractical, The judge model plays a pivotal role in curating effective attack prompt candidates for verification and update the attack algorithm in the next cycle.

Early attack frameworks such as GCG (greedy coordinate gradient) \cite{gcg2023} worked on accessible or locally deployed victim models based on the premise that their discovered attack prompts can be transferable to similarly deployed models in real applications, or be used to construct heuristics in black-box frameworks \cite{chaos2024}.
 They typically operate with a single generation cycle where the prompt creation process follows optimization algorithms incorporating evaluation results from previous prompts and their responses, and the judging mechanism is more primitive, based on detecting refusal patterns in the generated responses.
 Starting with a refusal response, the optimization algorithm modifies tokens in the attack prompt with the goal of suppressing victim's refusal behavior over iterations.
 In such methods, the optimization gradients are derived directly from the victim model, e.g., token-level likelihood.
 Most models generate relatively predictable refusal texts when encountering jailbreak request against their safety policies.
 There is ongoing research on refusal behavior detection and analysis in models \cite{xie2024sorrybench, extractcls2025}, which can also be leveraged by attackers to improve refusal-based judge components in these types of model-guided attacks, as well as to identify and fingerprint victim models in real-world applications \cite{pasquini2025llmmap}.

As models grow in sophistication and are integrated into applications beyond chatbots, model-guided attack frameworks for black-box settings have evolved alongside them to directly interact with various target applications.
 This requires more sophisticated judges in the attack loop\cite{pair2023}, similar to those employed in AI firewalls, to evaluate victim responses within application context and generate malicious prompts using model-guided heuristics rather than purely optimization-based methods.
 Unlike optimization-based attacks, these frameworks rely on iterative prompt refinement driven by attacker models.
 They enable more flexible exploration of the attack space without requiring internal model access and can infiltrate agentic AI ecosystems.
 The workflow in Algorithm \ref{algo:attack} also covers this type of model-guided attacks, which are designed similarly to well-known software fuzzing frameworks, as reflected in names such as GPTFuzz and PromptFuzz \cite{llmfuzzer2024, yu2023gptfuzzer, yu2024promptfuzz, gt2025}.
In these frameworks, the model-generated prompts in each cycle are tested against the victim application, and the resulting responses are evaluated by the judge model within the cycle.
Prompt-response pairs meeting the judge criteria of interest (no refusal, and potentially malicious behavior) are optionally verified and used to update the prompt creation heuristic and improve the success rate of the next generation cycle.
The prompt selection criteria based on the prompt-response judge score, represented by $\tau_g$ in Algorithm \ref{algo:attack}, can change dynamically per cycle in some frameworks.
During the early cycles probing for an effective attack heuristic, the attack prompts are expected to achieve marginal judge scores on the responses they evoke from the victim model.

\begin{algorithm}[t]
    \caption{Generic Model-Guided Attack Workflow}
    \label{algo:attack}
    \small
    \begin{algorithmic}[1]
        \Require victim model $V$ with defense mechanism function $F_D$; attacker judge model $J_A$; verifier $H$;
        generations $G$; attempts $N$; verification budget $K$;
        initial template set $\mathcal{S}_0$
        \Ensure verified set $\mathcal{D}$

        \State $\mathcal{D} \gets \emptyset$
        \State $\mathcal{S} \gets \mathcal{S}_0$

        \For{$g=1$ to $G$}
        \State $\mathcal{C}^N_g \gets \emptyset$ \Comment{candidate prompt set}
        \For{$i=1$ to $N$}
        \State $p_i \gets \textsc{CreatePrompt}(\mathcal{S}, \mathcal{C}^N_g)$
        \State $r_i \gets F_D(V(p_i))$  \Comment{filtered response}
        \State $s_i \gets J_A(p_i,r_i)~$ \Comment{effectiveness score}
        \State \textbf{if} $s_i \ge \tau_g$ \textbf{then} add $(p_i,r_i,s_i)$ to $\mathcal{C}^N_g$
        \EndFor

        \State $\mathcal{D}_g^{K} \gets \textsc{SelectCandidates}(\mathcal{C}^N_g, K)$
        \Comment{e.g., top-$K$}\vskip 4pt

        \State $\mathcal{D}_g^{K} \gets \textsc{Verify}(H,\mathcal{D}_g^{K})$
        \Comment{(optional for $g < G$)}

        \State $\mathcal{D} \gets \mathcal{D} \cup \mathcal{D}_g^{K}$
        \State $\mathcal{S} \gets \textsc{UpdateTemplate}(\mathcal{S}, \mathcal{C}^N_g, \mathcal{D})$
        \EndFor

        \State \Return $\mathcal{D}$
    \end{algorithmic}
    \normalsize
\end{algorithm}

Focusing on the inner cycle of the model-guided attack workflow in Algorithm \ref{algo:attack}, we consider a constrained verification setting where the attacker's success depends on having at least a verifiable effective candidate, i.e. $\mathcal{D}_g^K \ne \emptyset$, per generation cycle $g$.
 Our goal is to derive reasonable bounds on \textit{attacker's success rate} (ASR) in each cycle with respect to the detection performance of the defense system and the attacker's judge in the context of probabilistic binary hypothesis testing.

To analyze the effect of the defense system in either model-integrated or external form, we treat $V(p_i)$ as the unfiltered response that the victim model would produce for an attack prompt $p_i$.
The defense mechanism function $F_D$ then processes this response before it is delivered to the attacker.
If the response is judged harmless, $F_D$ passes it through; otherwise, under the detect-and-block strategy, it replaces the response with a clear refusal or deflection message.
As illustrated in Figure \ref{fig:fuzzing}, we consider the probabilities of the following complementary events concerning the victim model's response in our subsequent derivations:
\begin{table}[h]
    \centering
    \setlength{\tabcolsep}{6pt}
    \renewcommand{\arraystretch}{1.15}
    \begin{tabular}{lcc}
        \toprule
        \textbf{Victim's Response $V(p_i)$} & \textbf{Harmless} & \textbf{Harmful} \\
        \midrule
        \textbf{True Nature ($T$)} & $T_0$ & $T_1$ \\
        \textbf{Defense Judgment ($D$)}      & $D_0$ & $D_1$ \\
        \textbf{Attacker Judgment ($A$)}     & $A_0$ & $A_1$ \\
        \bottomrule
    \end{tabular}
\end{table}

Let $q \triangleq Pr(T_1)$ be the apriori (latent) probability of generating a harmful response from the underlying victim model, which is an unobserved constant quantity during the attack cycle. The performance of each detector is defined in terms of its \textit{type-I: false-positive (FP)} ($\alpha$) and \textit{type-II: false-negative (FN)} ($\beta$) error probabilities as follows:
\begin{align}\label{notationeq}
    \alpha_D \triangleq P(D_1 | T_0) ~,&~ \beta_D \triangleq P(D_0 | T_1)\\
    \alpha_A \triangleq P(A_1 | T_0) ~,&~ \beta_A \triangleq P(A_0 | T_1)
\end{align}
where $\alpha_D,\beta_D$ denote the defense system error probabilities and $\alpha_A,\beta_A$ denote the attacker's judge error probabilities.
Although both the attacker and the defense system aim to reduce these quantities, there is often an inherent trade-off between the type-I and type-II error probabilities. In particular, the defense system may operate under a Neyman-Pearson style criterion, where the false-positive rate $\alpha_D$ is constrained below a tolerable threshold while maximizing detection performance, since rejecting non-harmful responses is a costly error. Lastly, we consider the common logical case that the attacker judge functions independently from defense systems, for a given response of a known nature. In other words,  the attacker has no means of causally influencing the decision of defense system about an already generated response,
\begin{align}\label{indepeq}
    (A \perp\!\!\!\perp D) \mid T \;\Leftrightarrow\;  P(A,D | T) = P(A|T)~P(D|T)
\end{align}

We now evaluate the probability of event $A_1$, meaning that a received response is selected by the attacker's judge model in Algorithm \ref{algo:attack} as a successful candidate. We start by partitioning the event space using the complementary defense events $D_0$ and $D_1$:
\begin{align}\label{a1eq}
    P(A_1) = P(A_1 ,D_0) + P(A_1,D_1)
\end{align}
Here we assume the common refusal or deflection responses from the current detect-and-block defense strategy are properly detected as undesired by the attacker's judge model, i.e.
\begin{align}\label{rejeq}
    P(A_1, D_1) = 0.
\end{align}
We can further expand the expression in \eqref{a1eq} with respect to conditional probabilities about the true nature of the response ($T$) as follows:
\begin{align} \label{a1prob}
    P(A_1)  = \sum_{i=0}^1 P(A_1, D_0 \mid T_i)\,P(T_i)
\end{align}
Then we end up with the following expression using the independence condition in \eqref{indepeq} and the notations in \eqref{notationeq}:
\begin{align}
    P(A_1) = q\beta_D(1-\beta_A) + (1-q)(1-\alpha_D)\alpha_A
\end{align}

According to the workflow in Algorithm \ref{algo:attack}, an attacker's ultimate success depends on discovering true-positive attack prompts with verifiable effects. Given verification budget $K$, the attacker can optionally verify positive candidates selected in each cycle to avoid false-positive error propagation, or postpone verification until the last cycle. Irrespective of the verification schedule, we consider a cycle successful if it results in at least one verified true-positive attack prompt candidate. Hence, we use the following expression for the ASR per generation cycle. This can also serve as an upper-bound success criterion for frameworks such as \cite{autodan2023, attacktree}, where more than one true-positive candidate may be needed to properly update templates for future attack cycles.
\begin{align}\label{eq:defasr}
    \mathrm{ASR} &\triangleq P(\mathcal{D}_g^K\ne \emptyset) = 1 - P(\mathcal{D}_g^K = \emptyset)
\end{align}
We can then expand the probability term in \eqref{eq:defasr} with respect to the possible sizes of the candidate set, $|\mathcal{C}_g^N|$, using the total probability rule:
\begin{align}
    \mathrm{ASR} = 1-\sum_{n=0}^N P(\mathcal{D}_g^K = \emptyset~,~|\mathcal{C}^N_g| = n)
\end{align}
% Since having an empty candidate set implies no verified candidates neither, we can further adjust the range in sum above to
Since having an empty candidate set implies having no verified candidates, we can further adjust the range in the sum above to
\begin{align}\label{eq:defasr3}
    \mathrm{ASR} = P(\mathcal{C}^N_g \ne \emptyset) - \sum_{n=1}^N P(\mathcal{D}_g^K = \emptyset~,~|\mathcal{C}^N_g| = n)
\end{align}

The attack workflow in Algorithm \ref{algo:attack} permits progressive prompt generation that incorporates decision history about previous prompts. This may correlate judge decisions within the same cycle, so the attempts need not be i.i.d. However, once a candidate set is formed, the verification step evaluates selected prompt-response pairs independently and according to their true nature. We model the per-candidate verification success probability by $\eta_A$, the \textit{positive-predictive value (PPV)} of the attacker's judge, which is computed from the judge error probabilities using Bayes' theorem:
\begin{align}\label{eq:ppv}
    \eta_A&\triangleq P(T_1 | A_1)\\[8pt]\nonumber
    &=\frac{P(A_1~|~T_1)~P(T_1)}{P(A_1)}\\[8pt]\nonumber
    &=\frac{P(A_1, D_0~|~T_1)~P(T_1)}{P(A_1)}\\[8pt]\nonumber
    &=\frac{q\beta_D(1-\beta_A)}{q\beta_D(1-\beta_A) + (1-q)(1-\alpha_D)\alpha_A}\\[8pt]\nonumber
    &=\left(1 + \frac{1-q}{q}\frac{1-\alpha_D}{\beta_D}\frac{\alpha_A}{1-\beta_A}\right)^{-1}
\end{align}
Given the maximum verification budget $K$, we end up with the following for the sum terms in \eqref{eq:defasr3}
\begin{align}
    P(\mathcal{D}_g^K = \emptyset~,~|\mathcal{C}^N_g| = n) = P(|\mathcal{C}^N_g| = n)(1-\eta_A)^{\min(n,K)}.
\end{align}
We then use the following inequalities
\begin{align}
    (1-\eta_A)^K \le (1-\eta_A)^{\min(n,K)} \le (1-\eta_A)
\end{align}
to substitute the sum terms in \eqref{eq:defasr3} for $n \in \left[1,N\right]$ and deliver the following closed-form bounds on ASR:
\begin{align}\label{eq:asrbounds}
    \mathrm{ASR} &\le P(\mathcal{C}^N_g \ne \emptyset)(1 - (1-\eta_A)^K)\nonumber\\
    \mathrm{ASR} &\ge P(\mathcal{C}^N_g \ne \emptyset)(1 - (1-\eta_A))
\end{align}

The remaining term is $P(\mathcal{C}^N_g \ne \emptyset)$, the probability that at least one judge-positive candidate is selected in the cycle.
The i.i.d. prompt generation case is sufficient for this term to approach one as $N$ grows, but it is not necessary.
It can approach one as long as the search process remains non-degenerate, meaning that the conditional probability of discovering a judge-positive candidate does not vanish as the attacker continues searching.
This type of condition is standard in random-search convergence analyses \cite{solis1981minimization} and is discussed in Appendix \ref{app:nondegenerate-search}.
This particularly applies to intermediate attack cycles with more relaxed selection criteria (lower $\tau_g$) or the frameworks that continue the attempts until such candidate is found \cite{gcg2023,autodan2023}.
Under this less restrictive condition,
\begin{align}\label{eq:largeN}
    \lim_{N\to\infty}P(\mathcal{C}^N_g \ne \emptyset)=1.
\end{align}
The familiar expression $1-(1-P(A_1))^N$ is recovered as the homogeneous i.i.d. special case.

To derive the worst-case bound on ASR from the defense perspective, we evaluate the high-potency prompt regime $q=1$, corresponding to prompts that would trigger malicious responses from the underlying model whenever they are not blocked by the defense.
In this regime, the true-positive candidate probability is maximized, since
\begin{align}\label{eq:partial}
    \frac{\partial P(A_1)}{\partial q} > 0 , ~\frac{\partial \eta_A}{\partial q} > 0 \;\rightarrow\; \arg\max_q(\mathrm{ASR}) = 1.
\end{align}
This is a reasonable choice for the latent probability $q$ for worst-case attack prompts that can trigger malicious responses from the underlying model in the absence of any defense system.
In this worst-case scenario, the attacker can achieve $\eta_A=1$, implying that the automated judge in the cycle only needs to reliably detect refusals from the target model, which is achievable given the predictable refusal behavior of many current models.
Combining $\eta_A=1$ with \eqref{eq:largeN} yields
\begin{align}\label{eq:baseasr}
    \max_q(\mathrm{ASR}) = P(\mathcal{C}^N_g \ne \emptyset)
\end{align}
Under the homogeneous i.i.d. approximation, this large-$N$ behavior can be written in the familiar closed form
\begin{align}\label{eq:baseasr_iid}
    \max_q(\mathrm{ASR}) \approx 1 - (1-\beta_D(1-\beta_A))^N.
\end{align}
Thus, given $\beta_D(1-\beta_A)>0$, the attacker can increase its worst-case ASR arbitrarily close to $1$ by increasing the number of model-guided attempts per cycle $N$.

\begin{remark}[\textbf{Incompleteness of Defense Strategies}]
    Consistent with studies such as \cite{su2024mission_impossible}, which show the theoretical existence of prompts bypassing model alignment, the above analysis implies that detect-and-block defense strategies are incomplete against model-guided attacks capable of automated search in prompt space. Under a non-degenerate search process, their success probability can approach $1$ for a sufficiently large number of attempts $N$, provided that the effective pass-through probability $\beta_D(1-\beta_A)>0$.
    \begin{align}\label{eq:rem1}
        \lim_{N\to\infty} \max_q(\mathrm{ASR}) = 1 .
    \end{align}
\end{remark}

\section{Analysis of Misdirection as a Defense} \label{sec:propanalysis}
\begin{figure}[!t]
    \centering
    \includegraphics[width=\linewidth]{figs/proposed-llm-defense.png}
    \caption{Proposed detect-and-misdirect defense strategy against model-guided attacks}
    \label{fig:proposed}
\end{figure}
Our proposed defense strategy, depicted in Figure \ref{fig:proposed}, can be described as \textit{detect-and-misdirect}, complementing current defense strategies against the rising threat of model-guided automated attacks in agentic ecosystems.
Once it detects a pattern of iterative malicious prompts, the defense mechanism can employ misdirection techniques to produce engaging in-context responses instead of clear and predictable refusal messages.
The misdirection responses aim to impose a new category of false-positive errors on the attacker's automated judge model, i.e. \textit{misdirection-induced false-positives (MI-FP)}, whose probability $\gamma_A$ is defined using the notations introduced in Section \ref{sec:theory}:
\begin{align}
    \gamma_A \triangleq P(A_1|D_1)
\end{align}
This implies that the condition in \eqref{rejeq} no longer holds and
\begin{align}\label{eq:pa1}
    P(A_1) = \gamma_A P(D_1) + \sum_{i=0}^1 P(A_1, D_0 \mid T_i)\,P(T_i)
\end{align}
increases compared to \eqref{a1prob}, where
\begin{align}
    P(D_1) = q(1-\beta_D)+(1-q)\alpha_D .
\end{align}
The misdirection responses are only generated for prompts deemed malicious by the defense system ($D_1$ event) and have a non-functional true nature for the attacker ($T_0$).
Thus, it degrades the PPV of the candidate prompts discovered in the automated attack workflow in Algorithm \ref{algo:attack} by increasing the denominator of its expressions in \eqref{eq:ppv} while the numerator quantity remains intact.
\begin{align}\label{eq:ppv2}
    &\eta_A = \\\nonumber
    &\left(1 + \frac{(1-q)(1-\alpha_D)\alpha_A}{q\beta_D(1-\beta_A)} + \gamma_A\frac{q(1-\beta_D)+(1-q)\alpha_D}{q\beta_D(1-\beta_A)}\right)^{-1}
\end{align}
Significant reduction in $\eta_A$ defeats the purpose of model-guided iterative search, thus forcing the attacker to resort to exhaustive, and less practical, approaches to verify each prompt.
Since the conditions in \eqref{eq:partial} still hold regarding the a priori probability $q$, to derive the worst-case ASR we evaluate the quantities in \eqref{eq:ppv2} and \eqref{eq:pa1} as follows:
\begin{align}
    &P(A_1) = \beta_D(1-\beta_A) + \gamma_A(1-\beta_D) \label{eq:pa1q1}\\[8pt]
    &\eta_A =  \left(1+\frac{\gamma_A(1-\beta_D)}{\beta_D(1-\beta_A)}\right)^{-1} \label{eq:ppvq1} ~\mathrm{given}~~ q = 1
\end{align}

These expressions can be inserted into the ASR upper-bound in \eqref{eq:asrbound} and then extended for large $N$ according to \eqref{eq:largeN} to obtain the general bound in \eqref{uppereq}.
Nevertheless, we can derive the following tractable approximation for ASR under the homogeneous i.i.d. prompt generation condition with verification budget $K=1$ which we can use for numerical evaluation:
\begin{align}\label{eq:asrmax}
    &\max_{K=1}(\mathrm{ASR}) = \\\nonumber
    &\beta_D(1-\beta_A)\sum_{n=0}^{N-1}(\beta_A\beta_D+(1-\gamma_A)(1-\beta_D))^n.
\end{align}
This can be used to derive the following tighter worst-case ASR bounds:
\begin{align}\label{eq:asrbound}
    \beta_D(1-\beta_A)\le \max_{K=1}(\mathrm{ASR}) < \left(1+\frac{\gamma_A(1-\beta_D)}{\beta_D(1-\beta_A)}\right)^{-1},
\end{align}
confirming that the proposed strategy prevents the ASR convergence to $1$ shown in \eqref{eq:rem1} for large numbers of automated attempts $N$, provided that $\gamma_A > 0$.
%The plots in Figure \ref{fig:asr_dual} demonstrate how even small values of misdirection-induced false-positive rates ($\gamma \le 0.1$) curtail the attacker success gains from increasing automated iterations, particularly when the defense mechanism false-negative error rate is lower.

%\begin{remark}[\textbf{Efficacy \& Complementarity of LLMaaD}]
\begin{remark}[\textbf{Misdirection Strategy Impact}]
    \label{rem3}
    The following reduced upper-bound on the attacker success rate expresses the core strength of the detect-and-misdirect strategy against emerging model-guided automated attack frameworks.
    Integrated with existing defense strategies, replacing predictable refusal messages with in-context misdirection can disrupt the feedback mechanism in such frameworks, forcing attackers to resort to less-practical search approaches with exhaustive and manual verification.
    These misdirection responses have minimal collateral impact, since they are only generated in response to threats detected by the existing defense mechanisms.
    \begin{align} \label{uppereq}
        \boxed{\lim_{N\to\infty} \max_{q,\beta_A}(\mathrm{ASR}) \le  1-\left(1 + \frac{\beta_D}{\gamma_A(1-\beta_D)}\right)^{-K}}
    \end{align}
    This bound is the consequence of a non-zero misdirection-induced false-positive rate $\gamma_A$, which prevents the asymptotic attack success probability from reaching unity.
\end{remark}

In a broader context, the proposed detect-and-misdirect strategy can be applied to mitigate other automated attacks, such as black-box fuzzing, in scenarios where the attacker's evaluation feedback is subject to false-positive errors.
For example, this technique may not be applicable in classic software security scenarios where the attackers can programmatically trigger and verify their desired responses from their target (e.g. a key) without relying on semantic evaluators.
In contrast, an agentic AI system can demonstrate the attacker's desired malicious behavior in various forms that require the use of AI/ML evaluators, which are prone to false-positive errors.
This distinction is central to model-guided attacks where detecting predictable refusal or deflection behavior is often easier than verifying whether an apparently cooperative response contains a functional malicious output.
Misdirection complicates the attacker's evaluation task by shifting it from detecting refusal to extracting functional information from responses.
It can leverage known weaknesses of LLM-based evaluators, which may rely on superficial cues such as tone, structure, and perceived intent rather than strict semantic correctness \cite{chen2024humans, wang2024fair}.

Attackers may seek to adapt existing model-guided frameworks against misdirection through costly means such as judge ensembling and extra calibration.
However, they still need to start the early cycles with more relaxed candidate selection criteria (e.g. lower $\tau_g$ in Algorithm \ref{algo:attack}) to identify weaknesses of the target system with respect to prompt templates that they can make more effective over the next cycles.
Since they are more prone to accepting false positives (i.e. larger MI-FP rate $\gamma_A$) in early cycles, the effect of the detect-and-misdirect defense strategy extends beyond just the per-cycle ASR reduction and can divert subsequent attack cycles toward ineffective trajectories and resource exhaustion.
Substantial reduction of $\gamma_A$ in early cycles may require stricter automated evaluation or increased human verification, weakening the automation advantage of model-guided attacks.
Recent studies further show that LLM-as-a-judge systems exhibit important reliability limitations even for advanced models and may still require expert grounding or human oversight in complex evaluation tasks \cite{szymanski2025limitations, chen2024humans}.
Consequently, while attackers may temper the effect of misdirection through stronger evaluators or ensembles, reliably filtering misdirection-induced false positives ultimately pushes model-guided attacks toward more costly human verification and exhaustive search.

The techniques to deceive adversaries, often referred to as cyber deception, have been studied and applied in other fields of cybersecurity.
However, they often involve the deployment of decoy artifacts such as honeytokens and honeypots to divert attacks or analyze them for offensive security research.
Similarly, the proposed misdirection strategy can make apparent successes less trustworthy to the attacker, since they must spend additional resources distinguishing true bypasses from controlled false engagements.
This uncertainty may serve as both an operational and psychological deterrent against automated prompt injection and jailbreak attacks, consistent with prior observations that belief in deception mechanisms can deter adversarial behavior \cite{cydec-naval}.

%\begin{figure}[!htb]
%    \centering
%    \includegraphics[width=9cm]{figs/asr-dual.png}
%    \caption{Per-cycle Maximum  Attacker Success Rate (ASR) in \eqref{eq:asrmax} as a function of misdirection-induced false-positive rate $\gamma_A$ and number of model-guided iterations $N_A$, both in logarithmic scale, plotted for for $\beta_A = 0$ and $\beta_D \in \{0.05, 0.005\}$.}
%    \label{fig:asr_dual}
%\end{figure}

\section{Contextual Misdirection via Progressive Engagement Against Jailbreak Attacks}
\label{sec:cmpe}

To provide a proof of concept for the proposed detect-and-misdirect strategy, we present a conversational misdirection mechanism termed \textit{Contextual Misdirection via Progressive Engagement (CMPE)}, which we evaluate against automated jailbreak attacks.
The objective of the CMPE method is to substitute predictable refusal responses with responses that appear cooperative and semantically plausible, while subtly disrupting the attacker’s automated evaluation loop without violating the model's policies.

Illustrated in Algorithm \ref{algo:cmpe}, the CMPE method consists of three components.
First, a \textit{positive-intent preamble} establishes cooperative framing, which biases automated evaluators toward interpreting the response as desired by the attacker.
Second, a \textit{context expansion stage} reshapes the original prompt through token-level transformations (e.g., shuffling and lexical injection) and expands it into a coherent but non-operational narrative.
This stage preserves semantic proximity to the original query while avoiding execution of the malicious intent.
Third, a \textit{follow-up question} introduces conversational branching to engage the attacker and obscure the evasive intent of the defense mechanism.


Looking at the example in Figure \ref{fig:cmpe}, CMPE shares similarities with techniques studied in computational persuasion, where responses are designed to influence beliefs in interacting parties \cite{computational_persuasion}.
However, when used as a defense mechanism, CMPE is designed to disrupt adversarial evaluation mechanisms in automated jailbreak attacks by introducing controlled misdirection.
In particular, the generated responses exploit known limitations of LLM-based judges, which can rely on heuristic cues such as tone, structure, and perceived intent rather than strict semantic correctness \cite{zheng2023judging}.


The results in Section \ref{sec:num} demonstrate how CMPE can significantly increase misdirection-induced false positives in leading evaluation models and disrupt the automated jailbreak frameworks presented in \cite{pair2023, llmfuzzer2024}.
There in Appendix \ref{ap:ycmpe} we show that CMPE is safer and more effective than instructing an abliterated (refusal-suppressed) model to generate misdirection responses.
The CMPE method can also be used to create training datasets for integrating the detect-and-misdirect strategy as a learned technique in agentic security components.

\begin{algorithm}[t]
    \caption{Example of LLMaaD against jailbreak attacks:\\Contextual Misdirection via Progressive Engagement (CMPE)}
    \label{algo:cmpe}
    \small
    \begin{algorithmic}[1]
        \Require prompt $p$; lightweight model $M$; expansion budget $N_e$
        \Ensure misdirection response $r = s_{\text{pre}} \,\|\, s_{\text{ctx}} \,\|\, q$

        \State \textbf{(1) Positive-intent preamble}
        \State $s_{\text{pre}} \gets M(\text{make\_positive\_intent\_preamble},\, p)$
        \Comment{1--2 sentences}

        \State \textbf{(2) Safe context via prompt-reshaping}
        \State $w \gets \textsc{Tokens}(\textsc{StripPunct}(p))$
        \State $w \gets \textsc{InjectGlue}(w)$
        \State $w \gets \textsc{Shuffle}(w)$
        \State $\tilde{p} \gets \textsc{Join}(w)$
        \State $s_{\text{ctx}} \gets M(\text{expand\_text},\, \tilde{p},\, N_e)$
        \Comment{expand to $\approx N_e$ words}

        \State \textbf{(3) Follow-up question}
        \State $q \gets M(\text{make\_detailed\_followup\_question},\, p)$
        \Comment{1--2 sentences}

        \State $r \gets s_{\text{pre}} \,\|\, s_{\text{ctx}} \,\|\, q$
        \State \Return $r$
    \end{algorithmic}
    \normalsize
\end{algorithm}

\section{Experiments \& Numerical Results} \label{sec:num}
In this section \tclr{we} evaluate the efficacy of CMPE, our example misdirection method, in increasing \tclr{the} false-positive rate in leading safety evaluation models used in automated jailbreak attacks. We estimate the baseline maximum ASR for the attack scenario illustrated in Figure \ref{fig:fuzzing} where an AI firewall uses these models to detect-and-block malicious \tclr{prompt-response} pairs triggered by an automated jailbreak attack. We then show the reduction in the estimated maximum ASR for the comparable scenario in Figure \ref{fig:proposed} where the AI firewall uses \tclr{the CMPE-based detect-and-misdirect method} in response to similar \tclr{prompt-response} pairs.

\subsection{Jailbreak Prompts Dataset}
We evaluate our framework using a subset of 500 high-risk jailbreak prompts drawn from the \textbf{AdvBench dataset} introduced in the GCG work \cite{gcg2023}. These prompts represent a diverse set of illicit and harmful intent queries commonly used in red-teaming aligned language models. The dataset is publicly available\footnote{\url{https://github.com/llm-attacks/llm-attacks}} and has been widely adopted for benchmarking jailbreak attacks.

\subsection{Response Generation Model}\label{sec:resp-gen}
We use the \textbf{NeuralDaredevil-8B-abliterated} model\footnote{\url{https://huggingface.co/mlabonne/NeuralDaredevil-8B-abliterated}}, which is a lightweight \tclr{abliterated} model based on \tclr{the Meta LLaMA-3 architecture}, as the primary response generation engine due to its lightweight \tclr{nature} and strong benchmark performance for reduced alignment tasks. This model is used for two purposes:

\subsubsection{Malicious response generation}
For each jailbreak prompt in the AdvBench dataset, we generate candidate responses to estimate false-negative errors of judge models. To ensure quality of the samples, we apply a polling procedure with human oversight: each response is evaluated five times across six judge models in Section \ref{sec:judgemodels} with \tclr{scores normalized and averaged per judge}. We obtained 453 samples that achieved an average score above 0.75 on at least four out of six judges. The remaining 47 samples were then manually modified and retested to meet this criterion.

\subsubsection{Misdirection generation}
Since aligned models typically refuse to generate misleading content, we use the same model to construct misdirection responses. Specifically, for generating \tclr{500 CMPE responses} to each prompt according to Algorithm \ref{algo:cmpe}, we used this model for producing positive-intent preambles (line no. 2), context expansion (line no. 8 with $N_e = 200$ words), and follow-up questions (line no. 10).

\subsection{Evaluator Models}\label{sec:judgemodels}
For systematic evaluation of both malicious and misdirection responses under realistic model-guided attack conditions, we use a diverse set of modern automated judge models and omit simple keyword-based refusal detectors (e.g., used in early GCG work), as they are trivially bypassed by CMPE misdirection. We use the terms \emph{evaluator} and \emph{judge} interchangeably to denote models that assess response harmfulness, including both LLM-based judges and classifier-based detectors. When necessary, we explicitly distinguish LLM-based evaluators as \emph{LLM judges}.

\begin{itemize}

    \item \textbf{StrongREJECT (SR)} \cite{sr2024}: A \tclr{rubric-based instruction} to construct a LLM judge producing integer scores between $1$ and $5$ that measure the relevancy of responses to prompts and assign $0$ to rejection responses.

    \item \textbf{PAIR} \cite{pair2023}: An \tclr{instruction-based LLM judge} used in iterative jailbreak frameworks, producing integer scores between $1$ and $10$ reflecting response compliance with the prompted request and harmfulness.

    \item \textbf{HB-FT-LLaMA2-13B (HarmBench)} \cite{harmbench}: A fine-tuned LLaMA2 model that performs binary classification for harmfulness.

    \item \textbf{GPTFuzz-RoBERTa} \cite{yu2023gptfuzzer, llmfuzzer2024}:
    A RoBERTa-based classifier producing binary decisions along with confidence scores
    $p \in [0,1]$, representing the estimated probability of malicious content.
    We directly use $p$ as a continuous normalized harmfulness score.
\end{itemize}

\subsection{Estimating Judge Error Rates} \label{sec:err}
We evaluate each jailbreak prompt together with its corresponding
simulated malicious response generated as described in
Section~\ref{sec:resp-gen} across all judge models.
To account for stochastic variability in model outputs, each
prompt--response pair is evaluated $10$ times per judge, resulting
in a total of $5000$ scores per model.
The distribution of these scores is shown in Fig.~\ref{fig:vplot-malscores}.
For each judge, the harmfulness scores are normalized to the
$[0,1]$ range based on the empirical minimum and maximum scores,
enabling consistent cross-model comparison.
\begin{figure}[!htb]
    \centering
    \includegraphics[width=\linewidth]{figs/vplot-malscores.png}
    \caption{Per-judge normalized harmfulness score distributions for jailbreak response samples.}
    \label{fig:vplot-malscores}
\end{figure}

\begin{figure}[!htb]
    \centering
    \includegraphics[width=\linewidth]{figs/vplot-scores.png}
    \caption{Per-judge normalized harmfulness score distributions for CMPE misdirection response samples.}
    \label{fig:vplot-scores}
\end{figure}
We then repeat the same procedure using the same prompts paired with their corresponding misdirection responses generated via the CMPE method described in Section \ref{sec:resp-gen}, yielding another set of $5000$ harmfulness scores per judge. The resulting distributions are shown in Figure \ref{fig:vplot-scores}. These results illustrate that the proposed misdirection responses, while semantically safe, consistently receive high harmfulness scores across judges for a large fraction of jailbreak prompts.
\begin{table}[!htb]
    \centering
    \setlength{\tabcolsep}{6pt}
    \renewcommand{\arraystretch}{1.15}
    \label{tabbg}
    \begin{tabular}{lcc}
        \toprule
        \textbf{Judge / Evaluator} & $\hat{\beta}$ (FN) & $\hat{\gamma}$ (MI-FP) \\
        \midrule
        \textbf{SR-OSS-120B}       & \basecell{0.00883}   & \defcell{0.6276} \\
        \textbf{SR-Nemotron-30B}   & \basecell{0.004783}  & \defcell{0.7535} \\
        \textbf{PAIR-OSS-120B}     & \basecell{0.022075}  & \defcell{0.8524} \\
        \textbf{PAIR-Nemotron-30B} & \basecell{0.038852}  & \defcell{0.5667} \\
        \textbf{HB-FT-LLaMA2-13B}  & \basecell{0.060338}  & \defcell{0.4247} \\
        \textbf{GPTFuzz-RoBERTa-0.4B} & \basecell{0.211238}  & \defcell{0.7122} \\
        \bottomrule
    \end{tabular}
    \caption{Estimated average judge error rates.}
\end{table}
Given $10$ per-judge normalized scores for each response, we apply a decision threshold of $0.6$ to estimate the probability that a response is classified as harmful by each judge. For jailbreak responses, this quantity corresponds to the per-sample false-negative (FN) rate of the defender (i.e., harmful responses that are not blocked). For misdirection responses, the same quantity corresponds to the misdirection-induced (MI) false-positive (FP) rate on the attacker judge, since these responses are benign but intentionally constructed to appear harmful.

\begin{table*}[h]
    \centering
    \setlength{\tabcolsep}{4pt}
    \renewcommand{\arraystretch}{1.2}

    \begin{tabular}{lcccccc}
        \toprule
        \textbf{Defender $\downarrow$ / Attacker $\rightarrow$}
        & \textbf{SR-OSS-120B} & \textbf{SR-Nemotron-30B}
        & \textbf{PAIR-OSS-120B} & \textbf{PAIR-Nemotron-30B}
        & \textbf{HB-FT-LLaMA2-13B} & \textbf{GPTFuzz-RoBERTa} \\
        \midrule
        \textbf{SR-OSS-120B}
        & \basecell{0.2791} & \basecell{0.5811}
        & \basecell{0.4911} & \basecell{0.4797}
        & \basecell{0.3210} & \basecell{0.5275} \\

        \textbf{SR-Nemotron-30B}
        & \basecell{0.2484} & \basecell{0.2689}
        & \basecell{0.3809} & \basecell{0.2553}
        & \basecell{0.2553} & \basecell{0.3296} \\

        \textbf{PAIR-OSS-120B}
        & \basecell{0.8409} & \basecell{0.8927}
        & \basecell{0.2193} & \basecell{0.7065}
        & \basecell{0.5596} & \basecell{0.6948} \\

        \textbf{PAIR-Nemotron-30B}
        & \basecell{0.9722} & \basecell{0.9780}
        & \basecell{0.9494} & \basecell{0.6848}
        & \basecell{0.9386} & \basecell{0.9084} \\

        \textbf{HB-FT-LLaMA2-13B}
        & \basecell{0.9960} & \basecell{0.9976}
        & \basecell{0.9914} & \basecell{0.9931}
        & \basecell{0.1475} & \basecell{0.9524} \\

        \bottomrule
    \end{tabular}
    \caption{Baseline ASR under the detect-and-block strategy for $N=100$ iterations (higher values indicate greater vulnerability). Several attacker–defender configurations yield ASR values close to $1$, demonstrating the effectiveness of automated model-guided attacks in the absence of misdirection.}
    \label{tab:asr_base}
\end{table*}

\begin{table*}[h]
    \centering
    \setlength{\tabcolsep}{4pt}
    \renewcommand{\arraystretch}{1.2}

    \begin{tabular}{lcccccc}
        \toprule
        \textbf{Defender $\downarrow$ / Attacker $\rightarrow$}
        & \textbf{SR-OSS-120B} & \textbf{SR-Nemotron-30B}
        & \textbf{PAIR-OSS-120B} & \textbf{PAIR-Nemotron-30B}
        & \textbf{HB-FT-LLaMA2-13B} & \textbf{GPTFuzz-RoBERTa} \\
        \midrule
        \textbf{SR-OSS-120B}
        & \defcell{0.0014632} & \defcell{0.0066657}
        & \defcell{0.0038950} & \defcell{0.0055280}
        & \defcell{0.0028974} & \defcell{0.0055228} \\

        \textbf{SR-Nemotron-30B}
        & \defcell{0.0011319} & \defcell{0.0011161}
        & \defcell{0.0021360} & \defcell{0.0013287}
        & \defcell{0.0017510} & \defcell{0.0018432} \\

        \textbf{PAIR-OSS-120B}
        & \defcell{0.024098} & \defcell{0.025889}
        & \defcell{0.00064705} & \defcell{0.015184}
        & \defcell{0.010801} & \defcell{0.011434} \\

        \textbf{PAIR-Nemotron-30B}
        & \defcell{0.053730} & \defcell{0.048044}
        & \defcell{0.032878} & \defcell{0.014099}
        & \defcell{0.058840} & \defcell{0.029828} \\

        \textbf{HB-FT-LLaMA2-13B}
        & \defcell{0.083046} & \defcell{0.075957}
        & \defcell{0.054531} & \defcell{0.083207}
        & \defcell{0.00059035} & \defcell{0.040197} \\

        \bottomrule
    \end{tabular}
    \caption{ASR under the CMPE-based detect-and-misdirect strategy for $N=100$ iterations (lower values indicate stronger defense). The substantial reduction across all configurations confirms the effectiveness of misdirection-induced false positives in limiting attacker success.}
    \label{tab:asr_def}
\end{table*}

\subsection{\todo{ASR Evaluation \& Comparison (review)}}
\label{sec:asr_eval}

Using the per-sample error rates estimated in Section \ref{sec:err} for each judge, we compute the maximum ASR for each pair of attacker and defender models according to the formulation in Section \ref{sec:theory}. Specifically, we evaluate all pairwise combinations of attacker and defense evaluators to capture heterogeneous interactions in realistic model-guided attack scenarios.

The results are shown in Tables \ref{tab:asr_base} and \ref{tab:asr_def} for $N=100$ attacker iterations, averaged across all $500$ jailbreak samples. Table \ref{tab:asr_base} reports the baseline ASR under the detect-and-block strategy, while Table \ref{tab:asr_def} reports the ASR under the proposed CMPE-based detect-and-misdirect strategy.

From Table \ref{tab:asr_base}, we observe that the attacker ASR is consistently high across most attacker–defender combinations, with several configurations approaching values close to $1$. This is particularly evident for moderately aligned defense models, where automated search reliably discovers successful jailbreak prompts within $N=100$ iterations. These results empirically validate the theoretical observation in Section \ref{sec:theory} that, under detect-and-block defenses, ASR can approach unity as the number of automated attempts increases.

In contrast, Table \ref{tab:asr_def} shows that introducing CMPE-based misdirection responses significantly reduces the attacker ASR across all configurations. In many cases, the ASR drops by one to two orders of magnitude, with several entries falling below $10^{-2}$. This reduction is consistent across both strong and weak evaluator combinations, demonstrating the robustness of the proposed strategy.

The reduction in ASR can be directly attributed to the increase in misdirection-induced false-positive rates $\gamma_A$, which degrade the attacker's positive predictive value (PPV) and disrupt the effectiveness of iterative search. Notably, even configurations that exhibit near-perfect ASR in the baseline case show substantial degradation when CMPE is applied.

These results highlight a fundamental shift in the asymptotic behavior of model-guided attacks: while detect-and-block strategies allow the attacker to converge to high success rates, the detect-and-misdirect strategy enforces a bounded success regime, as predicted by the analytical results in Section \ref{sec:propanalysis}.

\subsection{End-to-End Evaluation Against Automated Attack Frameworks}
\label{sec:e2e_eval}

We conduct end-to-end adversarial evaluations using two state-of-the-art, model-guided attack frameworks: \textbf{GPTFuzz}~\cite{yu2023gptfuzzer,llmfuzzer2024} and \textbf{PAIR}~\cite{pair2023}. These experiments measure the operational attack success rate (ASR) when the full detect-and-block and detect-and-misdirect defense strategies are deployed against realistic automated attack campaigns.

\subsubsection{Experimental Configuration}

\mbox{}\\[8pt]
Table~\ref{tab:e2e_config} summarizes the models and roles used across all experiments. Each attack framework was evaluated on 50 adversarial goals independently sampled from the AdvBench dataset~\cite{gcg2023} to match each framework's original evaluation protocol; the subsets are not shared across frameworks. Both configurations were evaluated under the detect-and-block and detect-and-misdirect defense strategies with a maximum of 50 iterations allocated per prompt.

The choice of victim models is deliberate: we target the worst-case scenario for defense evaluation. Vicuna-13b exhibited the highest baseline ASR among aligned models in both the GPTFuzz and PAIR original papers, making it a representative upper bound for safety-aligned targets. The NeuralDaredevil-8B-abliterated model, having undergone refusal suppression, readily complies with jailbreak instructions without internal safety resistance. Together, these models span the spectrum from weakly aligned to entirely unaligned, providing a rigorous test of the defense framework's model-agnostic properties.

\begin{table}[t]
    \centering
    \small
    \setlength{\tabcolsep}{4pt}
    \renewcommand{\arraystretch}{1.15}
    \begin{tabular}{>{\raggedright\arraybackslash}p{2.0cm}
                    >{\raggedright\arraybackslash}p{3.8cm}
                    >{\raggedright\arraybackslash}p{1.8cm}}
        \toprule
        \textbf{Component} & \textbf{Model} & \textbf{Role} \\
        \midrule
        Victim (Target) & lmsys/\newline vicuna-13b-v1.5 & Primary \\
        Victim (Target) & mlabonne/\newline NeuralDaredevil-\newline 8B-abliterated & Secondary \\
        Mutator / Attacker & gpt-3.5-turbo & Prompt generation \\
        Attacker Judge & GPTFuzz-RoBERTa~\cite{yu2023gptfuzzer} & GPTFuzz Judging \\
        Attacker Judge & GPT-OSS-120B~\cite{pair2023} & PAIR Judging \\
        Defender & Llama-Guard-3-8B & Detection \\
        Reshaper & NeuralDaredevil-\newline 8B-abliterated & CMPE \\
        Secondary Judge & Claude Sonnet 4.6 & Validation \\
        \bottomrule
    \end{tabular}
    \caption{End-to-end evaluation model configuration.}
    \label{tab:e2e_config}
\end{table}

The \textbf{detect-and-block} strategy routes each victim model response through the Llama-Guard-3-8B defender model, serving as a firewall; responses classified as harmful are replaced with a refusal message. The \textbf{detect-and-misdirect} strategy instead invokes the CMPE reshaping algorithm (Algorithm~\ref{algo:cmpe}) to generate a misdirected response that appears to satisfy the attacker's query while containing no genuinely harmful content.

\subsubsection{Post-Hoc Validation via Secondary LLM Judge}

\mbox{}\\[8pt]
To ensure accurate evaluation of defense efficacy, we manually evaluated the final response produced by each framework for each of the 50 prompts and employed Claude Sonnet 4.6 as a secondary LLM judge to mitigate evaluation bias. Each jailbreak response was classified into one of two categories:
\begin{itemize}[leftmargin=*]
    \item \textbf{Jailbroken:} Genuinely harmful, actionable content that could cause real harm.
    \item \textbf{False Positive:} Safe, incoherent, or off-topic response incorrectly flagged as jailbroken by the attacker's judge.
\end{itemize}

The \textit{Atk JBs} column records the raw count of responses the attacker's judge accepted as successful jailbreaks. The \textit{True JBs}, \textit{False Pos}, and \textit{False Neg} columns are all derived from \textit{Atk JBs} through this manual validation process: \textit{True JBs} are attacker-accepted responses confirmed as genuinely harmful; \textit{False Pos} are attacker-accepted responses found to be safe, incoherent, or otherwise non-harmful (wrong or off-topic outputs incorrectly flagged as jailbreaks); and \textit{False Neg} are harmful responses that the attacker's judge failed to claim as jailbreaks, i.e., genuine harms missed by the attack framework's own judge. The \textit{Misdirected} column counts responses where the attacker accepted a reshaped output as a jailbreak, contributing to false positives under the detect-and-misdirect strategy. The \textit{Defended} column reports prompts where the attack framework exhausted its maximum iteration budget without claiming success.

\subsubsection{Results and Interpretation}

\mbox{}\\[8pt]
Tables~\ref{tab:gptfuzz_e2e} and~\ref{tab:pair_e2e} present the end-to-end results for GPTFuzz and PAIR, respectively.

\begin{table*}[t]
    \centering
    \setlength{\tabcolsep}{6pt}
    \renewcommand{\arraystretch}{1.2}
    \resizebox{\textwidth}{!}{%
    \begin{tabular}{llcccccccc}
        \toprule
        \textbf{Model} & \textbf{Defense} & \textbf{Total} & \textbf{Atk JBs} & \textbf{Misdirected} & \textbf{Defended} & \textbf{False Pos} & \textbf{False Neg} & \textbf{True JBs} & \textbf{Avg. Iter.} \\
        \midrule
        Vicuna      & Block     & 50 & 23 &  0 & 27 & 13 & 0 & 10 & 37.5 \\
        Vicuna      & Misdirect & 50 &  4 & 44 &  2 &  4 & 0 &  0 &  6.8 \\
        \midrule
        Abliterated & Block     & 50 & 22 &  0 & 28 & 12 & 2 &  8 & 34.0 \\
        Abliterated & Misdirect & 50 &  6 & 42 &  2 &  2 & 3 &  1 &  6.4 \\
        \bottomrule
    \end{tabular}}
    \caption{GPTFuzz end-to-end results (50 prompts, max 50 iterations, rigorously validated ASR).}
    \label{tab:gptfuzz_e2e}
\end{table*}

\begin{table*}[t]
    \centering
    \setlength{\tabcolsep}{6pt}
    \renewcommand{\arraystretch}{1.2}
    \resizebox{\textwidth}{!}{%
    \begin{tabular}{llcccccccc}
        \toprule
        \textbf{Model} & \textbf{Defense} & \textbf{Total} & \textbf{Atk JBs} & \textbf{Misdirected} & \textbf{Defended} & \textbf{False Pos} & \textbf{False Neg} & \textbf{True JBs} & \textbf{Avg. Iter.} \\
        \midrule
        Vicuna      & Block     & 50 &  5 &  0 & 45 &  2 & 0 &  3 & 46.4 \\
        Vicuna      & Misdirect & 50 &  3 & 36 & 11 &  1 & 0 &  2 & 22.8 \\
        \midrule
        Abliterated & Block     & 50 &  9 &  0 & 41 &  6 & 2 &  1 & 46.0 \\
        Abliterated & Misdirect & 50 &  3 & 46 &  1 &  1 & 0 &  2 & 13.2 \\
        \bottomrule
    \end{tabular}}
    \caption{PAIR end-to-end results (50 prompts, max 50 iterations, rigorously validated ASR).}
    \label{tab:pair_e2e}
\end{table*}

Under blocking, both frameworks achieved comparable validated true ASR: GPTFuzz reached 10/50 (20\%) on Vicuna-13b and 8/50 (16\%) on the abliterated model (avg.\ 34.0--37.5 iterations, 27--28 prompts defended), while PAIR achieved 3/50 (6\%) on Vicuna and 1/50 (2\%) on the abliterated model with fewer iterations (46.0--46.4) and a higher defended count (41--45), consistent with its more deliberate refinement strategy.

The detect-and-misdirect strategy produced markedly stronger defense: for GPTFuzz the true ASR collapsed to 0/50 (0\%) on Vicuna and 1/50 (2\%) on the abliterated model, with 44 and 42 responses misdirected respectively and convergence accelerating to just 6.4--6.8 iterations; for PAIR, the true ASR dropped to 2/50 (4\%) on both targets with convergence at 13.2--22.8 iterations and false positives rising to 95--96\%. In both cases, misdirection caused the attacker's judge to accept reshaped responses as genuine jailbreaks and terminate early, while the few remaining validated jailbreaks correspond exclusively to cases where Llama-Guard-3-8B failed to intercept the response, directly corresponding to the PPV reduction formalized in Eq.~\eqref{eq:ppv2} and the bounded ASR regime predicted in Section~\ref{sec:propanalysis}.

Post-hoc validation revealed that the RoBERTa attacker judge exhibits a false positive rate exceeding 50\% on block-mode responses, underscoring the necessity of secondary validation. Comparable misdirection performance across both the safety-aligned Vicuna-13b and the refusal-suppressed abliterated-8B confirms that misdirection framework proposed operates at the response layer, independent of the victim model's inherent alignment.


\subsection{Limitations}
\label{sec:limitations}

The scope of this work leaves several important considerations for future analysis and deployment.
The probabilistic derivation adopts a per-cycle analytical abstraction with a constrained verification budget.
This abstraction is useful for exposing the per-cycle scaling behavior of model-guided attacks, but it does not fully model how misdirection propagates across multiple dependent attack cycles, where false-positive candidates can influence future prompt templates and attacker heuristics.
In practice, this multi-cycle effect may strengthen the practical impact of misdirection, but it requires a more detailed sequential model of attacker adaptation.

Our empirical evaluation focuses primarily on jailbreak-style benchmarks and well-known attack frameworks.
Although the proposed mechanism is motivated by broader automated attacks in agentic AI systems, including prompt injection, further evaluation is needed in full tool-use and multi-agent environments.
We also did not conduct a dedicated study of the collateral effects of misdirection on benign human users.
The strategy may be particularly suitable for fully agentic environments where suspicious interactions can be handled without directly confusing a human user.
In chatbot settings, however, the frequency and conditions under which misdirection should be triggered depend on the confidence and policy judgment of the defending system.
Clear refusal messages may still be appropriate, and sometimes necessary, as warnings or policy explanations for benign users.


\subsection{Future Work}
\label{sec:future-work}

Future work should study how detect-and-misdirect can be integrated into
realistic defense architectures.
In this paper, CMPE is implemented as a separate response-generation mechanism activated after malicious behavior is
detected.
A natural next step is to investigate whether misdirection can be incorporated into model-level safety tuning, security reinforcement, or AI
firewall policies.
However, directly teaching large general-purpose models when
to misdirect may be difficult and could increase the risk of unintended
hallucination or confusing responses.
A more controlled deployment option is to
use a separate lightweight model or policy module that generates either refusal
or misdirection responses depending on the confidence, context, and severity of
the detected threat.

An important direction is to study detect-and-misdirect defenses in richer agentic
settings where the target system interacts with tools, memory, external
documents, and other agents.
Such environments create more complex attack
surfaces, but also provide additional opportunities for controlled response
shaping.

Finally, future work should evaluate the effect of misdirection in human-facing
security studies, such as controlled CTF-style experiments.
Such studies could
measure whether misdirection changes attacker behavior, confidence, and
resource allocation in non-automated or semi-automated scenarios.
This would
also clarify the relationship between detect-and-misdirect and broader cyber
deception mechanisms, including whether the belief that a target system may
deploy misdirection has a deterrent effect similar to the deception effects
observed in prior cybersecurity studies \cite{cydec-naval}.



\appendices
\section{Non-Degenerate Search Condition}
\label{app:nondegenerate-search}

The limit in \eqref{eq:largeN} does not require the prompt attempts inside a
generation cycle to be independent and identically distributed. It only
requires that the search process maintains a non-vanishing conditional
probability of producing a judge-positive candidate as attempts continue. This
type of non-degeneracy condition is standard in random-search convergence
analysis \cite{solis1981minimization}.

Let $A_{1,i}$ denote the event that the $i$-th attempt in a generation cycle is
selected as a candidate by the attacker's judge, and let $\mathcal{H}_{i-1}$
denote the history of previous attempts in the same cycle. Suppose that before
the first candidate is found, there exists some $\epsilon>0$ such that
\begin{align}
    P(A_{1,i}\mid \mathcal{H}_{i-1},
    A_{1,1}^{c},\ldots,A_{1,i-1}^{c}) \ge \epsilon
\end{align}
for all attempts $i$. Then the probability of producing no candidate after
$N$ attempts satisfies
\begin{align}
    P(\mathcal{C}_g^N=\emptyset)
    &=
    \prod_{i=1}^{N}
    P(A_{1,i}^{c}\mid \mathcal{H}_{i-1},
    A_{1,1}^{c},\ldots,A_{1,i-1}^{c})\\
    &\le (1-\epsilon)^N .
\end{align}
Therefore,
\begin{align}
    \lim_{N\to\infty}P(\mathcal{C}_g^N=\emptyset)=0,
    \qquad
    \lim_{N\to\infty}P(\mathcal{C}_g^N\ne\emptyset)=1 .
\end{align}
This proves that the convergence in \eqref{eq:largeN} holds under a less restrictive
condition than i.i.d. prompt attempts. The homogeneous i.i.d. expression
$1-(1-P(A_1))^N$ is recovered as a special case when each attempt has the same
candidate-selection probability and attempts are independent.

\begin{figure}[htb]
    \centering
    \includegraphics[width=\linewidth]{figs/spearman.png}
    \caption{xxxx}
    \label{fig:spearman}
\end{figure}


\section{CMPE Alternatives} \label{ap:ycmpe}
We have tested two other alternatives to generate misdirection responses for the prompts in the jailbreak dataset:
(1) directly prompting the \tclr{abliterated} model in Section \ref{sec:resp-gen} to generate persuasive, relevant misdirection responses, and
(2) \tclr{using} a simplified version of CMPE without generating the follow-up question (i.e., skipping step no. 10 in Algorithm \ref{algo:cmpe}).
\begin{figure}[htb]
    \centering
    \includegraphics[width=\linewidth]{figs/vplot-algocomp.png}
    \caption{Comparing CMPE with alternative misdirection methods for jailbreak attack scenario using GPT-OSS-120B configured as a StrongREJECT LLM judge for model-guided attack.}
    \label{fig:vplotalgos}
\end{figure}
The plot in Figure \ref{fig:vplotalgos} demonstrates the score distribution for CMPE in comparison to these two alternative methods using GPT-OSS-120B configured as a StrongREJECT LLM judge. It shows \tclr{that the abliterated misdirection responses achieve lower scores with higher variation}.
We have \tclr{investigated} some of the samples achieving high scores and they appear to be true harmful responses, which renders this method a risky approach to misdirection. The abliterated model sometimes ignores the intent of the instructions for safe misdirection and generates genuinely harmful responses, and when it complies with the instructions, the responses are often detected as irrelevant by the LLM judge.
This plot also highlights the effect of follow-up questions in misdirecting the LLM judge. While both variants remain safe and non-operational for attackers, there is a noticeable increase in the judge scores for CMPE responses compared to their counterparts without the follow-up question component.

\todo{The tighter and higher-scoring distribution of CMPE scores indicates improved effectiveness in inducing false-positive errors on attacker side.}



\bibliographystyle{IEEEtran}
\bibliography{llmaad-519}

\section*{LLM Usage Statement}
LLMs were used for language polishing and literature search during the
preparation of this manuscript. Claude Code was also used to assist with
implementation and deployment of proof-of-concept evaluation scripts.
All AI-assisted text modifications and generated code were reviewed, refined,
and validated by the authors. The theoretical derivations, technical
statements, concepts, experimental methodology, and conclusions presented
in this paper are original work of the authors, who take full responsibility
for the correctness, originality, and integrity of the document.

\end{document}



